\section*{第 \ref{ch:g12:stereochemistry} 章——立体化学：手性与异构体}

\begin{solution}{exo:g12:stereochemistry:1}
2-丁醇：碳 2（\ce{H}、\ce{OH}、\ce{CH3}、\ce{CH2CH3}）。
2-丙醇：无（碳 2 连有两个相同的 \ce{CH3}）。
2-氯丁烷：碳 2。乳酸：碳 2（\ce{H}、\ce{OH}、
\ce{CH3}、\ce{COOH}）。
\end{solution}

\begin{solution}{exo:g12:stereochemistry:2}
第一个：两个 \ce{Cl}（优先级高于 \ce{H}）位于同侧：Z。
第二个：左侧的 \ce{CH3} 优先级更高，位于上侧；右侧的
\ce{CH2CH3} 优先级更高，位于下侧：E。
\end{solution}

\begin{solution}{exo:g12:stereochemistry:3}
有\omterm{def:g12:stereochemistry:chiral}{手性}：手套、螺丝钉。无\omterm{def:g12:stereochemistry:chiral}{手性}：勺子、立方体（每个都有
一个镜面）。\ce{CHFClBr} 有\omterm{def:g12:stereochemistry:chiral}{手性}（碳\omterm{def:g7:atoms-and-molecules:atom}{原子}上连有四个不同的\omterm{def:g7:atoms-and-molecules:atom}{原子}）；
\ce{CH2Cl2} 无\omterm{def:g12:stereochemistry:chiral}{手性}。
\end{solution}

\begin{solution}{exo:g12:stereochemistry:4}
等量两种\omterm{def:g12:stereochemistry:enantiomer}{对映异构体}的\omterm{def:g6:pure-substances-and-mixtures:mixture}{混合物}。外消旋香芹酮含有
两种\omterm{def:g12:stereochemistry:enantiomer}{对映异构体}，因此鼻子同时接收两种气味。
\end{solution}

\begin{solution}{exo:g12:stereochemistry:5}
否：碳 1 连有两个相同的\omterm{def:g7:atoms-and-molecules:atom}{原子}，即 \ce{H} 和 \ce{H}；交换
它们所得仍为同一\omterm{def:g7:atoms-and-molecules:molecule}{分子}。
\end{solution}

\begin{solution}{exo:g12:stereochemistry:6}
\setchemfig{atom sep=2em}
\chemfig{C(-[2]CH_2CH_3)(-[4]H_3C)(<[:-30]OH)(<:[:-150]H)} \quad 镜像
\quad \chemfig{C(-[2]CH_2CH_3)(-[0]CH_3)(<[:-150]HO)(<:[:-30]H)}
\end{solution}

\begin{solution}{exo:g12:stereochemistry:7}
2-氯丁烷：2 种（一个\omterm{def:g12:stereochemistry:chiral}{手性碳原子}）。2,3-二氯丁烷：3 种（两个
\omterm{def:g12:stereochemistry:chiral}{手性碳原子}，两半相同：一对\omterm{def:g12:stereochemistry:enantiomer}{对映异构体}与一个内消旋
体）。2-戊烯：2 种（Z 和 E）。
\end{solution}

\begin{solution}{exo:g12:stereochemistry:8}
\omterm{def:g12:stereochemistry:enantiomer}{对映异构体}；\omterm{def:g12:stereochemistry:diastereomer}{非对映异构体}；\omterm{def:g12:stereochemistry:diastereomer}{非对映异构体}；同一\omterm{def:g7:atoms-and-molecules:molecule}{分子}（内消旋
酒石酸与其镜像可以重合）。
\end{solution}

\begin{solution}{exo:g12:stereochemistry:9}
在交错\omterm{def:g12:stereochemistry:conformation}{构象}中，两个 \ce{CH3} 基团位于相反方向，大基团
彼此尽可能远离：该\omterm{def:g12:stereochemistry:conformation}{构象}最稳定。在
重叠\omterm{def:g12:stereochemistry:conformation}{构象}中，两个 \ce{CH3} 相对，彼此拥挤。
\end{solution}

\begin{solution}{exo:g12:stereochemistry:10}
其两半各为 \ce{CH(OH)COOH}，二者互为镜像：一个
平面位于两个中心碳\omterm{def:g7:atoms-and-molecules:atom}{原子}之间，将\omterm{def:g7:atoms-and-molecules:molecule}{分子}切成两个互为镜像的
半部分。具有镜面的\omterm{def:g7:atoms-and-molecules:molecule}{分子}与其镜像相同：
无\omterm{def:g12:stereochemistry:chiral}{手性}。
\end{solution}

\begin{solution}{exo:g12:stereochemistry:11}
\setchemfig{atom sep=1.7em}
(Z)-2-戊烯 \chemfig{C(-[:150]H_3C)(-[:210]H)=C(-[:30]CH_2-CH_3)(-[:-30]H)}
和 (E)-2-戊烯 \chemfig{C(-[:150]H_3C)(-[:210]H)=C(-[:30]H)(-[:-30]CH_2-CH_3)}。
\end{solution}

\begin{solution}{exo:g12:stereochemistry:12}
与酒石酸类似，只是用 \ce{Br} 代替 \ce{OH}，\ce{CH3}
代替 \ce{COOH}：两个 \ce{Br} 均为实楔形，均为虚楔形（一对
\omterm{def:g12:stereochemistry:enantiomer}{对映异构体}），以及一实一虚，后者有镜面：即
内消旋体。共三种\omterm{def:g12:stereochemistry:stereoisomer}{立体异构体}。
\end{solution}

\begin{solution}{exo:g12:stereochemistry:13}
一半。化学家可以拆分两种\omterm{def:g12:stereochemistry:enantiomer}{对映异构体}（如巴斯德所做的，用的是
其他方法），或者只\omterm{def:g10:synthesis-yield:synthesis}{合成}有用的\omterm{def:g12:stereochemistry:enantiomer}{对映异构体}，采用
本身具有\omterm{def:g12:stereochemistry:chiral}{手性}的反应（\omterm{def:g12:stereochemistry:chiral}{手性}催化剂、酶或\omterm{def:g12:stereochemistry:chiral}{手性}起始
原料）。
\end{solution}

\begin{solution}{exo:g12:stereochemistry:14}
三种：(2E,4E)、(2Z,4Z) 和 (2E,4Z)，其中第三种与从另一端读取的
(2Z,4E) 为同一\omterm{def:g7:atoms-and-molecules:molecule}{分子}。
\end{solution}

\begin{solution}{exo:g12:stereochemistry:15}
$2^3 = 8$ 种\omterm{def:g12:stereochemistry:stereoisomer}{立体异构体}，即 4 对\omterm{def:g12:stereochemistry:enantiomer}{对映异构体}。每种
\omterm{def:g12:stereochemistry:stereoisomer}{立体异构体}有一个\omterm{def:g12:stereochemistry:enantiomer}{对映异构体}，并有 $8 - 2 = 6$ 个\omterm{def:g12:stereochemistry:diastereomer}{非对映异构体}。
\end{solution}

\begin{solution}{pb:g12:stereochemistry:1}
\textbf{1.} \ce{C4H6O6}；$M = 4 \times 12.0 + 6 \times 1.0 + 6 \times 16.0
= \qty{150.0}{g/mol}$。

\textbf{2.} 两个羧基（酸）和两个羟基
（醇）。

\textbf{3.} 两个中心碳\omterm{def:g7:atoms-and-molecules:atom}{原子}，各与 \ce{H}、\ce{OH}、
\ce{COOH} 和 \ce{CH(OH)COOH} 相连。

\textbf{4.} $2^2 = 4$。

\textbf{5.} 两者均为实楔形。

\textbf{6.} 两个 \ce{OH} 均为虚楔形；它不能与
L 重合：它是 D-酒石酸，即其\omterm{def:g12:stereochemistry:enantiomer}{对映异构体}。

\textbf{7.} 绕中心键旋转，使两个 \ce{OH} 朝向
同一方向，\omterm{def:g7:atoms-and-molecules:molecule}{分子}在两个中心碳\omterm{def:g7:atoms-and-molecules:atom}{原子}之间有一个平面，
该平面将一半反射到另一半上。

\textbf{8.} 否：它是内消旋酒石酸。

\textbf{9.} 在四种组合（实，实）、（虚，虚）、
（实，虚）、（虚，实）中，其中后两种为同一\omterm{def:g7:atoms-and-molecules:molecule}{分子}，
即内消旋形式，因为\omterm{def:g7:atoms-and-molecules:molecule}{分子}的两半相同。

\textbf{10.} L 与 D：\omterm{def:g12:stereochemistry:enantiomer}{对映异构体}。L 与内消旋体：\omterm{def:g12:stereochemistry:diastereomer}{非对映异构体}。

\textbf{11.} 在 \qty{169}{\celsius} 时同样如此：\omterm{def:g12:stereochemistry:enantiomer}{对映异构体}具有相同的
物理性质。

\textbf{12.} 否：它是 L 的\omterm{def:g12:stereochemistry:diastereomer}{非对映异构体}，是
另一种具有自身性质的化合物。

\textbf{13.} 鼻子中的受体具有\omterm{def:g12:stereochemistry:chiral}{手性}，能够区分
互为镜像的\omterm{def:g7:atoms-and-molecules:molecule}{分子}；温度计没有\omterm{def:g12:stereochemistry:chiral}{手性}，做不到这一点。

\textbf{14.} 否，这是两种化合物的\omterm{def:g6:pure-substances-and-mixtures:mixture}{混合物}；整体上不具有
光学活性，两种\omterm{def:g12:stereochemistry:enantiomer}{对映异构体}的旋光作用相互抵消。

\textbf{15.} 只含一种\omterm{def:g12:stereochemistry:enantiomer}{对映异构体}：每颗晶体都
仅由 L 或仅由 D 构成。

\textbf{16.} 各 \qty{5.0}{g}。

\textbf{17.} 内消旋体是另一种化合物，不属于
L 与 D 的\omterm{def:g12:stereochemistry:enantiomer}{外消旋混合物}。

\textbf{18.} 将每堆晶体分别\omterm{def:g2:mixing-and-dissolving:dissolve}{溶解}：两种溶液使偏振光
发生相同角度、方向相反的旋转（两种晶体
互为镜像）。

\textbf{19.} 共四种：其两个\omterm{def:g12:stereochemistry:chiral}{手性碳原子}所连基团不同
（一个连 \ce{Cl}，另一个连 \ce{OH}），因此没有任何组合具有对称面，
且任意两种组合都不是同一\omterm{def:g7:atoms-and-molecules:molecule}{分子}。

\textbf{20.} 三种：L、D 与内消旋体。
\end{solution}
